% !TeX root = ../机器学习笔记.tex
% 由主文档通过 \input 载入；请编译主文档。
\setcounter{chapter}{1} % 保留原书章号。
\chapter{方法比较与学习自检}
\label{app:ml-methods}

\section{从任务特征比较方法}

下表提供选择初始基线的思路，不宣称任何方法在某类任务上无条件最优。实际选择以合规的数据划分、验证表现及计算成本为依据。
\begin{longtable}{@{}L{3.0cm}L{5.0cm}L{5.7cm}@{}}
  \toprule
  方法族 & 主要假设或特点 & 重点检查 \\
  \midrule
  \endhead
  线性与逻辑回归 & 固定特征上的线性响应或线性分类分数 & 特征构造、共线性、尺度、正则化与概率校准。\\
  生成分类 & 类条件分布及先验有可用结构 & 分布设定、条件独立假设及少数类估计。\\
  近邻与核平滑 & 局部相似性包含预测信息 & 距离、维度、邻域覆盖与查询成本。\\
  树及其集成 & 递归切分或多模型组合 & 深度、叶样本量、调参预算与外推局限。\\
  核方法 & 核给出的特征几何适合任务 & 半正定性、尺度、惩罚强度及矩阵存储。\\
  聚类 & 距离、密度或连接结构反映有用分组 & 特征表示、簇形状、粒度与评价目标。\\
  PCA 与低秩近似 & 线性低维结构保留主要变化 & 中心化、尺度、重构与下游任务的差异。\\
  隐变量与混合模型 & 潜在成分有合理概率解释 & 可辨识性、初始化、退化与似然的界。\\
  神经网络 & 可训练表示适合任务结构 & 数据规模、梯度、训练诊断与泛化评价。\\
  序列和图模型 & 条件依赖或邻接结构有意义 & 推断成本、信息可用性与划分单位。\\
  \bottomrule
\end{longtable}

\section{逐章掌握程度与综合自检}

对每一方法，要求能独立回答五个问题：输入输出是什么，模型假设是什么，优化或推断目标是什么，算法怎样从初始化达到停止，结论如何在新数据上检验。进一步要求完成一个手算例题、一个维度与梯度检查、一个失败反例和一个公平的基线比较。能够调用软件只表明完成了计算操作，能够解释前提、误差来源与失败边界才表明理解了方法。





